\documentclass[10pt]{article}
\usepackage{../common/d3report}
\hypersetup{pdftitle={CS455 D3 Adversarial Test Report},pdfauthor={Mainak Sarkar}}
\newcommand{\RunDate}{2026-10-28}
\newcommand{\RunTotal}{107}
\newcommand{\RunPassed}{107}
\newcommand{\RunFailed}{0}
\newcommand{\RunFiles}{14}
\newcommand{\RunSeconds}{36.2}
\newcommand{\IntFiles}{10}
\newcommand{\IntTotal}{66}
\newcommand{\IntPassed}{66}
\newcommand{\BugTotal}{7}
\newcommand{\AdvFiles}{4}
\newcommand{\AdvTotal}{41}
\newcommand{\AdvPassed}{41}


\begin{document}

\makereporttitle{Adversarial Test Report}
  {Deliverable 3, Campus Ride-Pooling and Split-Fare System}
  {Prepared by Mainak Sarkar \quad Jira CS455-113, CS455-114 \quad Run of \RunDate}

\keybox{\textbf{Result.} \AdvTotal{} adversarial cases in \AdvFiles{} files try to break the app
through its real API routes against a real PostgreSQL database. After the fix below, all
\AdvPassed{} pass. The suite found one real weakness (an idempotency key replaying the wrong
request), fixed in PR \#87 with regression tests. It also confirmed that every race it staged is
serialized by the database, and it records one observation for the security report.}

\section{Purpose and method}
Ordinary tests show that the app does what it should when used as intended. Adversarial tests
use it as it should \emph{not} be used: at the wrong moment, with the wrong input, twice, out of
order, or several times at once. They run on the integration harness (see the integration test
report): the real route handlers, the real SQL, a real database emptied before every case.

\textbf{Pass criterion.} Every case must end in one of two ways: a clean refusal (a 4xx with a
reason code, nothing stored), or a consistent state (the rows obey the rules, for example never
a seat without an accepted request). A 500, a partial write or a misleading success fails the
case.

\textbf{Boundaries with the other D3 reports.} Access control (one user acting on another's
records), cross-site requests and secrets are the security report's subject; the 50-request
last-seat experiment (NFR-RD-05) is the concurrency report's. This suite keeps a small version of
each race as a regression test, run on every pull request.

\section{Attack categories}
\begin{tabularx}{\linewidth}{@{}>{\raggedright\arraybackslash}p{0.2\linewidth}X@{}}
\toprule
\textbf{Category} & \textbf{What is attacked, and why} \\
\midrule
Time boundaries & Every rule that depends on the clock, a moment either side of its boundary:
the lock an hour before departure (P-06), the five-minute no-show grace period (P-19), the
seven-day complaint window (FR-RD-13.1), the 72-hour rating window (P-16), the end of the
departure window. Off-by-one errors (\texttt{<} against \texttt{<=}) live here. \\
Input & Bodies that are not JSON objects, wrong types that a lenient parser might coerce
(\texttt{"600"}, 600.5, a timestamp number, month 13), oversized strings up to 1\,MB, malformed,
injection-looking and path-like identifiers, SQL and markup as text. \\
Replays & The same request sent again: a double-clicked accept, and an idempotency key sent
again with \emph{different} details. \\
Order & Actions that are valid in one state sent in another: requests and edits on a cancelled
ride, leaving or checking in after pickup started, pickup actions and SOS after completion,
acting again after leaving. \\
Races & Two to five requests at the same moment, each on its own database connection so they
really overlap: two accepts for the last seat, accept against withdraw, a five-fold
double-click, triple lock, five SOS taps, completion twice, the same offer twice; and the
per-address rate limit. \\
\bottomrule
\end{tabularx}

\section{Results}
\begin{center}\small
\captionof{table}{Adversarial test files, from the recorded run (\texttt{test-run/integration-results.json}).}
\begin{tabular}{@{}lrrr@{}}
\toprule
File (under \texttt{apps/web/test/integration/}) & Cases & Passed & Time (s) \\
\midrule
\tabinput{../test-run/generated/adversarial-files.tex}
\midrule
\textbf{Total} & \textbf{\AdvTotal} & \textbf{\AdvPassed} & \\
\bottomrule
\end{tabular}
\end{center}

Every case is listed in the appendix. The race file was also run five times in a row on its own,
passing each time, to check that its results do not depend on timing.

\section{Defect found: an idempotency key replayed the wrong request}\label{sec:bug}
\begin{tabularx}{\linewidth}{@{}>{\raggedright\arraybackslash}p{0.17\linewidth}X@{}}
\toprule
Case & \file{replays-and-state}: ``does not answer a request for one ride with a request on
another'', and ``refuses a different pickup sent with an idempotency key already used''. \\
Steps & A rider requests a seat on ride A with key \textit{K} (201). They then request a seat on
ride B, or on ride A with another pickup point, with the same key \textit{K}. \\
Expected & A refusal: \textit{K} belongs to a different request (FR-RD-06.9 makes a key replay
\emph{the same} request). \\
Actual & 200 \texttt{\{kind: "existing", requestId: <ride A's request>\}}. The page says ``This
request was already submitted''; no request exists on ride B, or the new pickup point is
silently dropped. The UI keeps its key across failed tries and can stay mounted across
client-side navigation, so a real rider could hit this. \\
Cause & \texttt{submitJoinRequest()} looked the key up and returned its request without
comparing what that request was for. \\
Fix & PR \#87: a key replays only the request it was first sent with (same ride and pickup
point); anything else gets 409 \texttt{idempotency\_key\_reused} and nothing is created. The
request form drops its key on that answer. \\
Regression & \file{bugs/idempotency-key-reuse} (3 tests; the two mismatch cases fail without the
fix) and 2 unit cases in \file{join-requests.test.ts}. Jira Bug, labels \texttt{d3-defect},
\texttt{found-by-adversarial}. \\
\bottomrule
\end{tabularx}

\section{What held, and why}
The other cases passed on the first run. For the races the reason is worth recording, since it is
what the concurrency report measures at scale:
\begin{itemize}
  \item \textbf{Last seat, accept against withdraw, double accept.} Each decision locks the
    rider's \texttt{users} row and the ride row (\texttt{FOR UPDATE} / \texttt{FOR NO KEY
    UPDATE}) and re-reads the seats inside that lock, so the second decision sees the first one's
    result: \texttt{ride\_full} or \texttt{request\_not\_active}, never an extra seat.
  \item \textbf{Double-clicked submit.} The rider's row lock serializes the five requests; the
    first stores the key, the other four find it and replay the same request.
  \item \textbf{SOS taps.} A partial unique index allows one active alert per person per ride,
    and the insert uses \texttt{ON CONFLICT DO NOTHING}, so taps cannot create a second alert even
    without a lock.
  \item \textbf{Completion twice, triple lock.} Trip steps lock the ride row and treat ``already
    there'' as done; the final shares are written once and still add up to the fare.
  \item \textbf{Input.} Route IDs are checked as UUIDs before any query; bodies go through one
    tolerant reader (\texttt{readJson}) and typed parsers that refuse rather than coerce; SQL is
    always parameterized, so hostile text is stored as text.
\end{itemize}

\section{Observation for the security report}
The per-address rate limit keys on the first value of the \texttt{X-Forwarded-For} header
(\texttt{clientIp()} in \file{src/lib/auth.ts}). The suite relies on this to give each call its own
address, and the rate-limit case shows that a new address starts a new count. Unless the server
sits behind a proxy that replaces that header, a client can choose its own address and avoid the
limit. Whether the deployment's proxy overwrites the header should be checked in the security
report; it is not counted as a defect here.

\section{Limitations}
\begin{itemize}
  \item The races use two to five requests; heavier contention is the concurrency report's
    experiment.
  \item Time boundaries are tested between six seconds and a minute either side, by moving data
    rather than the
    clock (see the integration report); exact-millisecond behaviour is covered by the unit tests
    of the rule functions.
  \item Access control and cross-user abuse are deliberately left to the security report.
\end{itemize}

\appendix
\section{Every adversarial case}
\small
\begin{longtable}{@{}>{\raggedright\arraybackslash}p{0.17\linewidth}>{\raggedright\arraybackslash}p{0.7\linewidth}l@{}}
\toprule
\textbf{File} & \textbf{Case (describe block: test)} & \textbf{Result} \\
\midrule
\endhead
\tabinput{../test-run/generated/adversarial-cases.tex}
\bottomrule
\end{longtable}

\end{document}
